%======================================================================
\section{Pair areas}\label{sec:pair}
%======================================================================

The generators $x$ and $z$ of $\GP$ act on the Pauli pair at the site $0$ only. To reach the pair at a site $u$ we conjugate by a word
$T_u$ that carries $0$ to $u$, a \emph{transporter} (\cref{fig:lamps}). Pairs at different sites commute in $\GP$, but an approximate
representation sees a relation only through a finite certificate: by \cref{lem:defect}, a word of area $k$ with respect to words of
defect at most $\eps$ has defect at most $k\eps$. \Cref{sec:sector} needs the $2^N$ pairs at the configurations of the first $N$ points of
ray~$1$ to commute up to small errors, so it needs their commutators to have small area with respect to one fixed finite set of words
$\RP$. This section proves that they have area at most $\Lambda_N\le2^{(\log N+23)^2}$ (\cref{thm:pair-area}), quasi-polynomial in $N$,
the logarithm of the number $2^N$ of pairs, and that the transported pairs square to the transported $J$ at linear cost (\cref{lem:J-transport}).

The proof is a computation with words in which every step involves at most three points of ray~$1$. Two lamps are first brought to
one translation word; the configuration that separates them is then cleared down to a single point by at most $N$ transvections, with the
translation word renormalized after each; what is left is the commutation of the lamps at two neighbouring sites, a word of $\RP$. Each
step is a \emph{local filling} (\cref{lem:local-filling}), whose cost is governed by the Dehn function $\dH$ of $H_3$ at scale $O(N)$: two
short words in $p,q$ that agree on the one or two points that matter differ by an element of a point stabilizer, which commutes with the
transported generator. This adapts~\cite[Theorem~5.1]{DM-lamps}, where the lamps are copies of $S_3$, to Pauli lamps.

\subsection{The certificate}\label{sec:pair-cert}

We use the notation of \cref{sec:group}. On the first ray we write $i$ for the point $(1,i)$, $e_i$ for $e_{(1,i)}$, $\tau_i$ for the
translation $\tau_{e_i}$, and $(i\ j)$ for the transposition of the points $i$ and $j$; thus $t=\tau_1$, $d=E_{2\leftarrow1}$ and
$s=[p,q]=(1\ 2)$ (\cref{group:gen}). By~\eqref{eq:affine-id} and \cref{def:GP}, a linear map fixes the site $0$ and so commutes with $x=x_0$
and $z=z_0$, an element $h\in A$ conjugates the pair at $v$ to the pair at $h(v)$, and $J$ is central (\cref{prop:J-central}). We write $J$
also for the word $xzxz$.

\begin{definition}[The certificate $\RP$]\label{def:RP}
Put
\[
s=[p,q],\qquad B=psp^{-1},\qquad t_m=p^{m-1}\,t\,p^{1-m}\quad(m=1,2,3).
\]
The \emph{certificate} $\RP$ consists of the $27$ words of \cref{tab:RP} in the letters $p,q,t,d,x,z$: thirteen \emph{host words} in
$p,q,t,d$ (families \textsf{H}, \textsf{Q}, \textsf{S}, \textsf{C}, \textsf{A}) and fourteen \emph{lamp words} (families \textsf{Z},
\textsf{L}, \textsf{P}, \textsf{K}). Expanded literally, the words have length at most $30$; we take each to be its free reduction, which
has length at most $22$. In particular $\RP$ contains $x^2$, $z^2$ and $(xz)^4$. In this section areas are taken with respect to $\RP$, and
we write $\Area=\Area_{\RP}$.
\end{definition}

\begin{table}[ht]
\centering\small
\begin{tabular}{@{}l p{0.38\textwidth} p{0.31\textwidth} r r@{}}
\toprule
 & words & meaning & number & longest\tabularnewline
\midrule
\textsf{H} & \raggedright $s^2$,\ $(sps\,p^{-1})^3$,\ $[s,p^2sp^{-2}]$,\ $psp^{-1}qs^{-1}q^{-1}$ & \raggedright the presentation $P_H$ of $H_3$ & 4 & 22\tabularnewline
\textsf{Q} & \raggedright $t^2$ & \raggedright translations are involutions & 1 & 2\tabularnewline
\textsf{S} & \raggedright $[t,sp]$, $[t,sq]$, $[d,Bsp]$, $[d,Bsq]$ & \raggedright stabilizers of prototypes commute with them & 4 & 16\tabularnewline
\textsf{C} & \raggedright $[t_1,t_2]$ & \raggedright translations commute & 1 & 8\tabularnewline
\textsf{A} & \raggedright $d\,t_1d^{-1}(t_1t_2)^{-1}$,\ $[d,t_2]$,\ $[d,t_3]$ & \raggedright how $d$ moves the translations at $1,2,3$ & 3 & 12\tabularnewline
\midrule
\textsf{Z} & \raggedright $[\ell,a]$ for $\ell\in\{p,q,d\}$, $a\in\{x,z\}$ & \raggedright linear maps fix the lamp at $0$ & 6 & 4\tabularnewline
\textsf{L} & \raggedright $x^2$,\ $z^2$,\ $(xz)^4$ & \raggedright the lamp at $0$ is $D_4$ & 3 & 8\tabularnewline
\textsf{P} & \raggedright $[a,tbt^{-1}]$ for $a,b\in\{x,z\}$ & \raggedright the lamps at $0$ and $e_1$ commute & 4 & 8\tabularnewline
\textsf{K} & \raggedright $[J,t]$, with $J=xzxz$ & \raggedright $J$ commutes with $t$ & 1 & 10\tabularnewline
\midrule
 & total & & 27 & \tabularnewline
\bottomrule
\end{tabular}
\caption{The certificate $\RP$: thirteen host words in $p,q,t,d$ and fourteen lamp words, written with the abbreviations of
\cref{def:RP}. The last column is the largest length of a word of the family after free reduction.}
\label{tab:RP}
\end{table}

\begin{lemma}\label{lem:RP-valid}
Every word of $\RP$ represents $1$ in $\GP$.
\end{lemma}

\begin{proof}
By \cref{group:gen}(a), $s=(1\ 2)$ and $B=(2\ 3)$. We check the families in turn.
\begin{itemize}[itemsep=0pt,leftmargin=2em]
\item[\textsf{H}] These are the relators of $P_H$, which hold in $H$ (\cref{sec:group-houghton}).
\item[\textsf{Q}] $t$ is a translation of the $\F_2$-space $V$, so $t^2=1$.
\item[\textsf{S}] $sp$ and $sq$ fix the point $1$, and $Bsp$ and $Bsq$ fix the points $1$ and $2$; by~\eqref{eq:affine-id} they commute
with $t=\tau_1$, respectively with $d=E_{2\leftarrow1}$.
\item[\textsf{C}] Translations commute.
\item[\textsf{A}] $t_m=\tau_m$ because $p^{m-1}$ maps the point $1$ to $m$, and $dt_md^{-1}=\tau_{d(e_m)}$ by~\eqref{eq:affine-id}, with
$d(e_1)=e_1+e_2$, $d(e_2)=e_2$ and $d(e_3)=e_3$.
\item[\textsf{Z}] $p$, $q$ and $d$ are linear, so they fix the site $0$ and commute with $x$ and $z$.
\item[\textsf{L}] $x$ and $z$ are involutions, and $(xz)^4=J^2=1$.
\item[\textsf{P}] $tbt^{-1}$ lies in the pair at the site $t(0)=e_1\ne0$, and pairs at distinct sites commute in $\EV$.
\item[\textsf{K}] $J$ is central.\qedhere
\end{itemize}
\end{proof}

The $27$ words are listed in the ancillary file \texttt{relators\_pauli.txt}, and a script checks them independently in a faithful model
of $\GP$ (\cref{app:cert}).

\paragraph{Areas.} Area is invariant under free equality, inversion and conjugation, and subadditive under products. If $\mu U\nu$ and
$\mu U'\nu$ are words, then $\mu U\nu=(\mu UU'^{-1}\mu^{-1})\,\mu U'\nu$ freely, so replacing a subword $U$ by $U'$ changes the area by at
most $\Area(UU'^{-1})$; we call this a \emph{replacement of cost $\Area(UU'^{-1})$}. Finally, $[w,v_1v_2]=[w,v_1]\cdot v_1[w,v_2]v_1^{-1}$
and $[w,v^{-1}]=v^{-1}[w,v]^{-1}v$ freely, so for all words $w,v$ and every factorization $v=v_1^{\pm1}\cdots v_m^{\pm1}$ into words $v_l$
\begin{equation}\label{eq:pair-comm}
\Area([w,v])=\Area([v,w])\le\sum_{l=1}^m\Area([w,v_l]).
\end{equation}
The family \textsf{H} is the presentation $P_H$ of~\eqref{eq:PH}, so a word $w$ in $p,q$ that represents $1$ in $H$ has
$\Area(w)\le\dH(|w|)$. The function $\dH$ is finite and nondecreasing (\cref{cor:dehn-qp}), and it is the only input about $H$ beyond
explicit computations with permutations.

\subsection{Short words in Houghton's group}\label{sec:pair-houghton}

All transport takes place among the first $n$ points of ray~$1$. For $n\ge1$ let $\Sigma_n\le H$ be the group of permutations supported
in $\{1,\dots,n\}$, identified with $\Sym(n)$. Recall that $\sigma_k=p^ksp^{-k}$ is the transposition $(k{+}1\ \,k{+}2)$
(\cref{group:gen}). Lengths of words in $p,q$ are counted with $s$ written out as $pqp^{-1}q^{-1}$.

\begin{lemma}[Finite permutations]\label{lem:pair-transp}
\leavevmode
\begin{enumerate}[label=(\alph*),itemsep=1pt]
\item For $1\le i<j$ the word $\omega_{i,j}=p^{i-1}(sp)^{j-i-1}(sp^{-1})^{j-i-1}s\,p^{1-i}$ represents $(i\ j)$ and has length
$10j-8i-8<10j$.
\item Every element of $\Sigma_n$ that moves $m\ge2$ points is represented by a word in $p,q$ of length less than $10(m-1)n$.
\item Let $y_1,\dots,y_l\in\{1,\dots,n\}$ be distinct. Put $\pi_0=1$ and, for $m=1,\dots,l$, $w_m=\pi_{m-1}^{-1}(y_m)$ and
$\pi_m=\pi_{m-1}(m\ w_m)$, where $(m\ m)=1$. Then $\pi[y_1,\dots,y_l]:=\pi_l$ maps $m$ to $y_m$ for $m\le l$, and it is a product of at
most $l$ transpositions of points of $\{1,\dots,n\}$. The product of the corresponding words $\omega_{m,w_m}$, also written
$\pi[y_1,\dots,y_l]$, has length less than $10ln$.
\end{enumerate}
\end{lemma}

\begin{proof}
(a) Conjugating $(i{+}1\ \,j)$ by $\sigma_{i-1}=(i\ \,i{+}1)$ gives $(i\ j)$ for $j>i+1$, so by induction on $j-i$
\[
(i\ j)=\sigma_{i-1}\sigma_i\cdots\sigma_{j-3}\;\sigma_{j-2}\;\sigma_{j-3}\cdots\sigma_i\sigma_{i-1}.
\]
In this product the powers of $p$ telescope: $\sigma_k\sigma_{k+1}=p^ksps\,p^{-k-1}$ and $\sigma_{k+1}\sigma_k=p^{k+1}sp^{-1}s\,p^{-k}$, and
multiplying out gives $\omega_{i,j}$. Its length is $2(i-1)+5(j-i-1)+5(j-i-1)+4=10j-8i-8$.

(b) A permutation that moves $m$ points is a product of at most $m-1$ transpositions of those points, each represented by (a) by a word
of length less than $10n$.

(c) By induction $\pi_{m-1}(k)=y_k$ for $k<m$. Since $y_m\ne y_k$ for $k<m$, the point $w_m=\pi_{m-1}^{-1}(y_m)$ is not among $1,\dots,m-1$,
so $w_m\ge m$ and $(m\ w_m)$ fixes $1,\dots,m-1$. Hence $\pi_m(k)=y_k$ for $k<m$, and $\pi_m(m)=\pi_{m-1}(w_m)=y_m$. Each $w_m$ lies in
$\{1,\dots,n\}$, because $\pi_{m-1}\in\Sigma_n$. There are at most $l$ nontrivial factors, each a word of length less than $10n$ by (a).
\end{proof}

That a point stabilizer of $H$ is isomorphic to $H$ is known~\cite[Example~3.6]{Cor06}. The next lemma gives such isomorphisms explicitly
on generators, with short words.

\begin{lemma}[Point stabilizers]\label{lem:pair-stab}
For $k\ge1$ let $H^{(k)}$ be the pointwise stabilizer in $H$ of the points $1,\dots,k$, put $\zeta_k=\sigma_{k-1}\cdots\sigma_1\sigma_0$ and
$K_k=(\zeta_kp,\ \zeta_kq)$.
\begin{enumerate}[label=(\alph*),itemsep=1pt]
\item There is an isomorphism $H\to H^{(k)}$ that sends $p,q$ to $\zeta_kp,\zeta_kq$ and the transposition $(i\ j)$ to
$(i{+}k\ \ j{+}k)$. In particular $K_k$ generates $H^{(k)}$.
\item Every element of $\Sigma_n\cap H^{(k)}$ that moves $m\ge2$ points is represented by a word with fewer than $10(m-1)n$ letters from
$K_k^{\pm1}$.
\end{enumerate}
Here $\zeta_1=s$ and $\zeta_2=Bs$, so $K_1=(sp,sq)$ and $K_2=(Bsp,Bsq)$ are the pairs of the family \textsf{S}; their letters have length
at most $5$ and $11$ in $p,q$.
\end{lemma}

\begin{proof}
(a) Let $\theta_k:Y\to Y\setminus\{1,\dots,k\}$ map $(1,j)$ to $(1,j+k)$ and fix rays $2$ and $3$. For $g\in H$ let $g^{[k]}$ act as
$\theta_kg\theta_k^{-1}$ on the image of $\theta_k$ and fix the points $1,\dots,k$. Then $g^{[k]}$ is eventually a translation on each ray,
and $g\mapsto g^{[k]}$ is an injective homomorphism $H\to H^{(k)}$. It is onto, because an element $h\in H^{(k)}$ restricts to a
permutation of the image of $\theta_k$, and $\theta_k^{-1}h\theta_k$ lies in $H$. Since $\theta_k$ maps the point $i$ to $i+k$, $(i\ j)^{[k]}=(i{+}k\ \ j{+}k)$. The element
$\zeta_k$ is the cycle $1\mapsto k{+}1\mapsto k\mapsto\dots\mapsto2\mapsto1$ of points of ray~$1$. Hence $\zeta_kp$ fixes the points
$1,\dots,k$, maps $(1,j)$ to $(1,j+1)$ for $j>k$, maps $(2,1)$ to $(1,k+1)$, and agrees with $p$ on the other points; that is,
$\zeta_kp=p^{[k]}$. Likewise $\zeta_kq=q^{[k]}$.

(b) An element $g\in\Sigma_n\cap H^{(k)}$ that moves $m$ points is $g'^{[k]}$ for some $g'\in\Sigma_{n-k}$ that moves $m$ points. Write $g'$
as a word in $p,q$ of length less than $10(m-1)(n-k)$ by \cref{lem:pair-transp}(b), and substitute $\zeta_kp,\zeta_kq$ for $p,q$.

Finally $\zeta_1=\sigma_0=s$ and $\zeta_2=\sigma_1\sigma_0=Bs$. With $s=pqp^{-1}q^{-1}$ and $B=psp^{-1}$, the words $sp$, $sq$ have length
$5$, and $Bsp$, $Bsq$ have length $11$.
\end{proof}

\subsection{Transporters and local fillings}\label{sec:pair-local}

\begin{definition}[Transporters]\label{def:transporter}
Let $n\ge3$. For distinct points $i,j\in\{1,\dots,n\}$ put
\[
U_i=\pi[i],\qquad \tau_i=U_i\,t\,U_i^{-1},\qquad U_{ij}=\pi[j,i],\qquad X_{ij}=U_{ij}\,d\,U_{ij}^{-1},
\]
words in $p,q,t,d$ built from the words of \cref{lem:pair-transp}(c). For $u\subseteq\{1,\dots,n\}$ with elements $i_1<\dots<i_m$ put
$T_u=\tau_{i_1}\cdots\tau_{i_m}$, with $T_\emptyset$ the empty word. We identify $u$ with the configuration $\sum_{i\in u}e_i\in V$.
\end{definition}

Since $U_i$ maps the point $1$ to $i$, the word $\tau_i$ represents the translation $\tau_i=\tau_{e_i}$ by~\eqref{eq:affine-id}, and we use
the same symbol for both. Since $U_{ij}$ maps $1$ to $j$ and $2$ to $i$, the word $X_{ij}$ represents the transvection $E_{i\leftarrow j}$.
Translations commute, so $T_u$ represents $\tau_u$, and $T_uxT_u^{-1}$ and $T_uzT_u^{-1}$ represent the Pauli pair $x_u,z_u$ at the site
$u$. By \cref{lem:pair-transp}, $|U_i|<10n$ and $|U_{ij}|<20n$, and $U_i$ is a product of at most one transposition and $U_{ij}$ of at most
two.

The elements $t$ and $d$ have \emph{prototypes} $P_t=\{1\}$ and $P_d=\{1,2\}$: by~\eqref{eq:affine-id}, an element of $H$ that fixes $P_f$
pointwise commutes with $f$. So two transporting words that agree on the prototype give the same conjugate. The next lemma prices this,
and the lemma after it lists the four local identities used in the proof of \cref{thm:pair-area}. They are the Pauli versions
of~\cite[Lemmas~4.1 and~4.2]{DM-lamps}. The transporting word may end with a power of $p$, because the words of the family \textsf{A}
use $t_2$ and $t_3$.

\begin{lemma}[Change of transporter]\label{lem:exchange}
Let $n\ge3$ and $f\in\{t,d\}$. Let $U$ and $W_0$ be words in $p^{\pm1},q^{\pm1}$ of length less than $30n$ that represent elements of
$\Sigma_n$, such that $W_0^{-1}U$ moves at most $10$ points. Let $W=W_0p^j$, where $j=0$ if $f=d$ and $j\in\{0,1,2\}$ if $f=t$. If the
permutations represented by $W$ and $U$ agree on $P_f$, then
\[
\Area\bigl(WfW^{-1}\cdot Uf^{-1}U^{-1}\bigr)\ \le\ \mathcal C(n):=2\,\dH(1500n)+130n .
\]
\end{lemma}

\begin{proof}
\emph{Step 1: a word for the stabilizer element.} Let $k=|P_f|$ and let $g$ be the element represented by $W^{-1}U$. It fixes $P_f$
pointwise, so $g\in H^{(k)}$. If $j=0$, then $g\in\Sigma_n$ moves at most $10$ points, and by \cref{lem:pair-stab}(b) it is represented by a
word $Z$ with fewer than $90n$ letters from $K_k^{\pm1}$. If $j\ge1$, then $f=t$, $k=1$, and $g=(sp)^{-j}g'$ with
$g'=(sp)^jp^{-j}\cdot W_0^{-1}U$. Here $(sp)p^{-1}=s$ and $(sp)^2p^{-2}=sB$ lie in $\Sigma_3$, so $g'\in\Sigma_n$ moves at most $13$ points,
and $g'$ fixes the point $1$ because $g$ and $sp$ do. By \cref{lem:pair-stab}(b), $g'$ is represented by a word $Z'$ with fewer than $120n$
letters from $K_1^{\pm1}$, and we put $Z=(sp)^{-j}Z'$. In both cases $Z$ has fewer than $122n$ letters from $K_k^{\pm1}$, each of length at
most $11$ in $p,q$.

\emph{Step 2: the filling.} The word $\Omega=W^{-1}UZ^{-1}$ in $p,q$ represents $1$ in $H$ and has length less than
$(30n+2)+30n+11\cdot122n\le1500n$, so $\Area(\Omega)\le\dH(1500n)$ by the family \textsf{H}. For each letter $\ell$ of $K_k$ the word
$[f,\ell]$ is in the family \textsf{S}, so $\Area([f,Z])<122n$ by~\eqref{eq:pair-comm}. Since $U=W\Omega Z$ freely,
\[
WfW^{-1}\cdot Uf^{-1}U^{-1}=W\,(f\Omega f^{-1})\,[f,Z]\,\Omega^{-1}\,W^{-1}
\]
freely, and its area is at most $2\dH(1500n)+122n\le\mathcal C(n)$.
\end{proof}

\begin{lemma}[Local fillings]\label{lem:local-filling}
Let $n\ge3$ and
\[
L(n)=10\,\dH(1500n)+700\,n .
\]
Let $i,j,k\in\{1,\dots,n\}$ with $i\ne j$, and $a,b\in\{x,z\}$. Each of the following words represents $1$ in $\GP$ and has area at most
$L(n)$.
\begin{enumerate}[label=(\alph*),itemsep=1pt]
\item $\tau_i^2$, of area at most $1$, and $[\tau_i,\tau_j]$, of area at most $4\mathcal C(n)+1$.
\item $X_{ij}\tau_kX_{ij}^{-1}\tau_k^{-1}$ if $k\ne j$, and $X_{ij}\tau_jX_{ij}^{-1}\tau_i^{-1}\tau_j^{-1}$; of area at most
$5\mathcal C(n)+1$.
\item $[X_{ij},a]$, of area at most $2|U_{ij}|+1<40n+1$.
\item $[a,\tau_ib\tau_i^{-1}]$, of area at most $4|U_i|+1<40n+1$.
\end{enumerate}
\end{lemma}

\begin{proof}
The words represent $1$: in (a) and (b) because translations are commuting involutions and $E_{i\leftarrow j}\tau_kE_{i\leftarrow
j}^{-1}=\tau_{E_{i\leftarrow j}(e_k)}$, with $E_{i\leftarrow j}(e_k)=e_k$ for $k\ne j$ and $E_{i\leftarrow j}(e_j)=e_i+e_j$; in (c) because
$E_{i\leftarrow j}$ is linear; in (d) because $\tau_ib\tau_i^{-1}$ lies in the pair at the site $e_i\ne0$. Since $5\mathcal C(n)+1\le L(n)$
and $40n+1\le L(n)$, it remains to prove the stated area bounds.

\emph{Changes of transporter.} In (a) and (b) we replace a subword $Uf^{\pm1}U^{-1}$ by $Wf^{\pm1}W^{-1}$ as in \cref{lem:exchange}; this
is a replacement of cost at most $\mathcal C(n)$, since the two words compared are conjugate to the word of the lemma or to its inverse.
Below, $W_0$ is a product of at most three transpositions, of length less than $30n$ (\cref{lem:pair-transp}(c)), and
$U\in\{U_i,U_j,U_k,U_{ij}\}$ is a product of at most two, so $W_0^{-1}U$ moves at most $10$ points.

(a) The word $\tau_i^2$ is freely $U_it^2U_i^{-1}$, a conjugate of the word $t^2$. For $[\tau_i,\tau_j]$ let $W_0=\pi[i,j]$, which maps $1$ to
$i$ and $2$ to $j$. Replace both occurrences of $\tau_i^{\pm1}$ by $(W_0tW_0^{-1})^{\pm1}$, and both of $\tau_j^{\pm1}$ by
$(W_0t_2W_0^{-1})^{\pm1}=(W_0p\,t\,p^{-1}W_0^{-1})^{\pm1}$: these are four changes of transporter, since $W_0$ and $U_i$ map $1$ to $i$, and
$W_0p$ and $U_j$ map $1$ to $j$. What is left is $W_0[t_1,t_2]W_0^{-1}$, a conjugate of the word of the family \textsf{C}. So the area is at
most $4\mathcal C(n)+1$.

(b) Let $W_0=\pi[j,i,k]$, or $W_0=\pi[j,i]$ if $k\in\{i,j\}$, so that $W_0$ maps $1$ to $j$ and $2$ to $i$, and let $m\in\{1,2,3\}$ be the
point with $W_0(m)=k$; thus $m=1$ exactly when $k=j$, and $m=2$ exactly when $k=i$. Replace $X_{ij}^{\pm1}$ by $(W_0dW_0^{-1})^{\pm1}$, a
change of transporter because $W_0$ and $U_{ij}$ agree on $P_d=\{1,2\}$; replace $\tau_k^{\pm1}$ by $(W_0t_mW_0^{-1})^{\pm1}$, where
$t_m=p^{m-1}tp^{1-m}$ and $W_0p^{m-1}$ maps $1$ to $k$; and, when $k=j$, replace $\tau_i^{\pm1}$ by $(W_0t_2W_0^{-1})^{\pm1}$. If $k\ne j$,
four changes turn the word into $W_0(dt_md^{-1}t_m^{-1})W_0^{-1}$ with $m\in\{2,3\}$, a conjugate of a word of the family \textsf{A}. If
$k=j$, five turn it into $W_0(dt_1d^{-1}t_2^{-1}t_1^{-1})W_0^{-1}$, again a conjugate of one. So the area is at most $5\mathcal C(n)+1$.

(c) The word $X_{ij}=U_{ij}dU_{ij}^{-1}$ has $2|U_{ij}|+1$ letters, each in $\{p,q,d\}^{\pm1}$, and $[\ell,a]$ is a word of the family
\textsf{Z} for $\ell\in\{p,q,d\}$. By~\eqref{eq:pair-comm}, $\Area([X_{ij},a])\le2|U_{ij}|+1<40n+1$.

(d) Freely the word is $U_i\,[U_i^{-1}aU_i,\ t(U_i^{-1}bU_i)t^{-1}]\,U_i^{-1}$. By~\eqref{eq:pair-comm} and the family \textsf{Z}, replacing
$U_i^{-1}a^{\pm1}U_i$ by $a^{\pm1}$ costs at most $\Area([U_i^{-1},a])\le|U_i|$, and likewise for $b$. Four such replacements leave
$U_i[a,tbt^{-1}]U_i^{-1}$, a conjugate of a word of the family \textsf{P}. So the area is at most $4|U_i|+1<40n+1$.
\end{proof}

\subsection{The pair-area theorem}\label{sec:pair-thm}

\begin{theorem}[Pair areas]\label{thm:pair-area}
Let $N\ge1$ and $n=\max(3,N)$. For every $u\subseteq\{1,\dots,N\}$ the transporter $T_u$ is a word in $p,q,t$ of length $|T_u|\le21Nn$ that
represents the translation $\tau_u$, so that $T_uxT_u^{-1}$ and $T_uzT_u^{-1}$ represent the Pauli pair at the site $u$. For distinct
$u,w\subseteq\{1,\dots,N\}$ and $a,b\in\{x,z\}$,
\[
\Area\bigl([T_uaT_u^{-1},\,T_wbT_w^{-1}]\bigr)\ \le\ \Lambda_N:=20N^2L(n),
\]
with $L(n)=10\dH(1500n)+700n$ as in \cref{lem:local-filling}. Moreover, with $C_0=4.5\cdot10^4$,
\[
14000N^2n\ \le\ \Lambda_N\ \le\ C_0N^2\bigl(\dH(C_0N)+N\bigr)\ \le\ 2^{(\log N+23)^2}.
\]
\end{theorem}

\paragraph{Overview of the proof.} In $\GP$ the statement is a short computation. Conjugation by $\tau_u$ turns the commutator of the
lamps at $u$ and $w$ into that of the lamps at $0$ and $u_0=u+w\ne0$. Fix a point $j$ of $u_0$. The transvection $E_{i\leftarrow j}$
toggles the point $i$ in every configuration that contains $j$, so the product $G$ of the at most $N-1$ transvections $E_{i\leftarrow j}$ over
the other points $i$ of $u_0$ is a linear map that carries $u_0$ to $e_j$ and fixes $0$. So the commutator is conjugate to that of the lamps
at $0$ and $e_j$, and these commute by the family \textsf{P}. The proof follows this computation word by word, and pays for each identity
it uses with a local filling (\cref{lem:local-filling}). \emph{Step~1} merges $T_u^{-1}T_w$ into the single translation word $T_{u_0}$ by
$O(N^2)$ swaps. \emph{Step~2} applies the transvections one at a time, and after each one renormalizes the translation word, so that it
never has more than $N+1$ factors; this costs $O(N)$ local fillings per transvection. \emph{Step~3} moves the lamp letters across $G$,
which is linear, and ends with one word of the family \textsf{P}. Each local filling costs at most $L(n)$, so the total is $O(N^2L(n))$.

\begin{example}\label{pair:ex}
Let $N=3$, $u=\{1,3\}$ and $w=\{2\}$, so that $T_u=\tau_1\tau_3$ (as in \cref{fig:lamps}) and $T_w=\tau_2$. Up to conjugation by $T_u$, the
commutator is $[a,ZbZ^{-1}]$ with $Z=T_u^{-1}T_w=\tau_3^{-1}\tau_1^{-1}\tau_2$. In Step~1, two words $\tau_i^2$ turn $Z$ into $\tau_3\tau_1\tau_2$,
and two swaps sort it into $T_{u_0}=\tau_1\tau_2\tau_3$, where $u_0=\{1,2,3\}$. In Step~2 we take the pivot $j=3$. Conjugation by $X_{13}$,
which represents $E_{1\leftarrow3}$, fixes $\tau_1$ and $\tau_2$ and turns $\tau_3$ into $\tau_3\tau_1$, up to local fillings:
\[
X_{13}\,\tau_1\tau_2\tau_3\,X_{13}^{-1}\ \leadsto\ \tau_1\tau_2\tau_3\tau_1\ \leadsto\ \tau_1\tau_1\tau_2\tau_3\ \leadsto\ \tau_2\tau_3=T_{\{2,3\}},
\]
by two swaps and one word $\tau_1^2$. Conjugation by $X_{23}$ turns $\tau_2\tau_3$ into $\tau_2\tau_3\tau_2$, then $\tau_2\tau_2\tau_3$, then
$\tau_3$. So $G=X_{23}X_{13}$ carries the site $u_0$ to $e_3$; indeed $E_{2\leftarrow3}E_{1\leftarrow3}(e_1+e_2+e_3)=e_3$. In Step~3 the
letters $a^{\pm1},b^{\pm1}$ are moved across $G$, which represents a linear map, and what is left is $[a,\tau_3b\tau_3^{-1}]$, a local filling
of type (d). Every identity used involves at most three points.
\end{example}

\begin{proof}[Proof of \cref{thm:pair-area}]
\emph{Step 0: the transporters.} Translations commute, so $T_u$ represents $\tau_u$. Each $\tau_i$ is a word in $p,q,t$ of length
$2|U_i|+1<20n+1$, so $|T_u|\le20Nn\le21Nn$. Write $L=L(n)$; by \cref{lem:local-filling} each of its words has area at most $L$, and
$L\ge700n$.

\emph{Step 1: one translation word.} Freely $[T_uaT_u^{-1},T_wbT_w^{-1}]=T_u[a,ZbZ^{-1}]T_u^{-1}$ with $Z=T_u^{-1}T_w$. Replace each of the
at most $N$ factors $\tau_i^{-1}$ of $T_u^{-1}$ by $\tau_i$, at cost $\Area(\tau_i^{-2})\le1$ each. Sort the resulting product of at most $2N$
factors $\tau_i$ into increasing order of $i$ by at most $\binom{2N}2$ swaps of adjacent distinct factors, each of cost at most $L$
(\cref{lem:local-filling}(a)), and cancel the at most $N$ adjacent equal pairs, at cost $1$ each. The result is $T_{u_0}$ with
$u_0=u+w\ne0$, so $\Area(ZT_{u_0}^{-1})\le2N+N(2N-1)L\le2N^2L$. Since $Z^{\pm1}$ occurs four times in $[a,ZbZ^{-1}]$,
\[
\Area\bigl([T_uaT_u^{-1},T_wbT_w^{-1}]\bigr)\le8N^2L+\Area\bigl([a,T_{u_0}bT_{u_0}^{-1}]\bigr).
\]

\emph{Step 2: normalized clearing.} Let $r\le N$ be the number of points of $u_0$, choose one of them, $j$, as the \emph{pivot}, and list the
others as $i_1,\dots,i_{r-1}$. Put $G_l=X_{i_lj}\cdots X_{i_1j}$ and $u_l=u_0+e_{i_1}+\dots+e_{i_l}$, which is $u_0$ with $i_1,\dots,i_l$
removed; so $u_{r-1}=\{j\}$. We claim
\[
\Area\bigl(G_lT_{u_0}G_l^{-1}\,T_{u_l}^{-1}\bigr)\le2lrL\qquad(0\le l\le r-1).
\]
For $l=0$ the word is freely trivial. Given the claim for $l$, put $i=i_{l+1}$ and $X=X_{ij}$. Freely
$G_{l+1}T_{u_0}G_{l+1}^{-1}=X(G_lT_{u_0}G_l^{-1})X^{-1}$, which becomes $XT_{u_l}X^{-1}$ at cost at most $2lrL$. Freely $XT_{u_l}X^{-1}$ is
the product, in order, of the words $X\tau_{i'}X^{-1}$ over the points $i'$ of $u_l$. By \cref{lem:local-filling}(b), at a total cost of at
most $rL$, each of them becomes $\tau_{i'}$, except that $X\tau_jX^{-1}$ becomes $\tau_j\tau_i$. The new factor $\tau_i$ has a partner
$\tau_i$ in $T_{u_l}$, since $i$ is a point of $u_l$; at most $r-1$ swaps make the two adjacent, and $\tau_i^2$ cancels them at cost $1$.
What remains is $T_{u_{l+1}}$, and this step costs at most $rL+(r-1)L+1\le2rL$, which proves the claim. The translation word is back in
normal form after every transvection and never has more than $r+1$ factors.

\emph{Step 3: a single point.} Let $G=G_{r-1}$ and $\Theta=GT_{u_0}G^{-1}\tau_j^{-1}$, of area at most $2r(r-1)L\le2N^2L$ by Step~2. Freely
$T_{u_0}=G^{-1}\Theta\tau_jG$. Substituting this into the four occurrences of $T_{u_0}^{\pm1}$ in $[a,T_{u_0}bT_{u_0}^{-1}]$ and deleting the
four occurrences of $\Theta^{\pm1}$ costs at most $8N^2L$, and leaves
\[
a\,G^{-1}\tau_jG\,b\,G^{-1}\tau_j^{-1}G\,a^{-1}\,G^{-1}\tau_jG\,b^{-1}\,G^{-1}\tau_j^{-1}G .
\]
The word $G$ is a product of $r-1$ words $X_{ij}$, so $\Area([G,c])\le(r-1)(40n+1)$ for $c\in\{x,z\}$ by~\eqref{eq:pair-comm} and
\cref{lem:local-filling}(c), and likewise for $[G,c^{-1}]=c^{-1}[G,c]^{-1}c$. Replacing $Gb^{\pm1}G^{-1}$ by $b^{\pm1}$, and then, after
conjugating by $G$, $Ga^{\pm1}G^{-1}$ by $a^{\pm1}$, costs at most $4(r-1)(40n+1)$ and leaves a conjugate of $[a,\tau_jb\tau_j^{-1}]$, of
area at most $40n+1$ by \cref{lem:local-filling}(d). Altogether, since $4N(40n+1)+40n+1\le205Nn\le N^2L$,
\[
\Area\bigl([a,T_{u_0}bT_{u_0}^{-1}]\bigr)\le8N^2L+4N(40n+1)+40n+1\le9N^2L,
\]
and with Step~1 the area of the commutator is at most $17N^2L\le\Lambda_N$.

\emph{Step 4: the size of $\Lambda_N$.} Since $L(n)\ge700n$, $\Lambda_N\ge14000N^2n$. Since $n\le3N$ and $\dH$ is nondecreasing,
$L(n)\le10\dH(4500N)+2100N$, so $\Lambda_N\le200N^2\dH(4500N)+42000N^3\le C_0N^2(\dH(C_0N)+N)$. By~\ref{H-exp},
$\dH(C_0N)\le2^{(\log(C_0N)+6)^2}$, and also $N\le2^{(\log(C_0N)+6)^2}$; so with $\ell=\log N\ge0$ and $\log C_0<15.46$,
\[
\log\Lambda_N\le\log(2C_0)+2\ell+(\ell+\log C_0+6)^2\le16.46+2\ell+(\ell+21.46)^2\le\ell^2+45\ell+477\le(\ell+23)^2 .
\qedhere
\]
\end{proof}

\begin{remark}[Growth of $\Lambda_N$]\label{rem:pair-growth}
By \cref{thm:pair-area}, $14000N^2n\le\Lambda_N\le2^{(\log N+23)^2}$ for all $N\ge1$. The factor $N^2$ counts local moves: at most $N$
transvections, each followed by at most $N$ swaps. The Dehn function of $H_3$ enters only through $L(n)$, at the scale $n=\max(3,N)$ of the
transport. So all $\binom{2^N}2$ pairs of lamps on $N$ points commute at area $O(N^2\dH(C_0N))$, which is $2^{O(\log^2N)}$ by~\ref{H-exp}.
Under~\ref{H-poly}, $\dH(m)\le C_Hm^{p_H}$ with $p_H\ge1$, and then
\[
\Lambda_N\le C_0N^2\bigl(C_H(C_0N)^{p_H}+N\bigr)\le C'N^{2+p_H}=C'N^{\kappa}\qquad(N\ge1),\qquad C'=C_0\bigl(C_HC_0^{p_H}+1\bigr),
\]
polynomial in $N$.
\end{remark}

\subsection{Transporting the central involution}\label{sec:pair-J}

The lower bounds compare the image of $J$ with the products $x_uz_ux_uz_u$ at the various sites. Transporting $J$ is cheap, and the
transported Pauli pair squares to the transported $J$ with no relation at all.

\begin{lemma}[Transporting $J$]\label{lem:J-transport}
Let $N\ge1$, $n=\max(3,N)$ and $u\subseteq\{1,\dots,N\}$. Then
\[
\Area\bigl(T_uJT_u^{-1}J^{-1}\bigr)\le4|T_u|\le84Nn,
\]
and $(T_uxT_u^{-1}\,T_uzT_u^{-1})^2=T_uJT_u^{-1}$ freely.
\end{lemma}

\begin{proof}
The word $T_uJT_u^{-1}J^{-1}$ is $[T_u,J]$, and the letters of $T_u$ are $p^{\pm1}$, $q^{\pm1}$ and $t^{\pm1}$. For $g\in\{p,q\}$,
\eqref{eq:pair-comm} applied to the four letters of $J=xzxz$ and the family \textsf{Z} give $\Area([g,J])\le4$. For the letter $t$,
$[t,J]=[J,t]^{-1}$ is the inverse of the word of the family \textsf{K}. By~\eqref{eq:pair-comm} applied to the letters of $T_u$,
$\Area([T_u,J])\le4|T_u|$, and $|T_u|\le21Nn$ by \cref{thm:pair-area}. The second claim holds because $T_u^{-1}T_u$ cancels freely, so the
word is freely $T_u(xz)^2T_u^{-1}=T_uJT_u^{-1}$.
\end{proof}
